\documentclass{article}
\usepackage[UTF8]{ctex}
\usepackage{amsmath,amssymb,amsfonts}
\usepackage{enumitem}
\usepackage{geometry}
\geometry{a4paper, margin=1in}
% === 错题本排版优化（不参与题目标识） ===
\setlength{\parindent}{0pt}
\setlength{\parskip}{0.35em plus 0.10em minus 0.08em}
\setlength{\emergencystretch}{3em}
\setlist{itemsep=0.25em, topsep=0.40em, parsep=0pt}
\allowdisplaybreaks[2]
% === 错题本排版优化结束 ===
\begin{document}
\section*{中国矿业大学2024年高等代数真题}
\section*{19．中国矿业大学2024年高等代数试卷}
一．填空题．每题4分，共40分．\\

\subsection{题目1}
1．若 1 是多项式 \(A x^{4}+B x^{2}+1\) 的二重根，则 \(A=\) \(\_\_\_\_\)，\(B=\) \(\_\_\_\_\)．\\
\(\_\_\_\_\)

\subsection{题目2}
2．设 \(A\) 是3阶方阵，\(|A|=-2\) ，把它按列分块成 \(\left(A_{1}, A_{2}, A_{3}\right)\) ，则方阵（ \(A_{3}-A_{1}, 3 A_{2}, A_{1}\) ）的行列式的值为 \(\_\_\_\_\)。

\subsection{题目3}
3．设 \(A=\left(\begin{array}{ccc}1 & 2 & -2 \\ 4 & k & 3 \\ 3 & -1 & 1\end{array}\right), B\) 为三阶非零矩阵，且 \(A B=O\) ，则 \(k=\) \(\_\_\_\_\)．\\
\(\_\_\_\_\)．

\begin{itemize}
  \item[A．] 交换 \(A^{*}\) 的第1列与第2列得到 \(B^{*}\) 。
  \item[B．] 交换 \(A^{*}\) 的第1行与第2行得到 \(B^{*}\)
  \item[C．] 交换 \(A^{*}\) 的第1行与第2行得到 \(-B^{*}\) 。
  \item[D．] 交换 \(A^{*}\) 的第1行与第2行得到 \(-B^{*}\) 。
\end{itemize}

\subsection{题目4}
5．设 \(A=\left(\alpha_{1}, \alpha_{2}, \alpha_{3}, \alpha_{4}\right)\) 为 4 阶方阵，其中 \(\alpha_{1}, \alpha_{2}, \alpha_{3}, \alpha_{4}\) 分别是 \(A\) 的列向量，且 \(|A|=0,|A|\) 的代数余子式 \(A_{12} \neq 0, A^{*}\) 是 \(A\) 的伴随矩阵，则方程组 \(A^{*} X=0\) 有一个基础解系是 \(\_\_\_\_\)．\\
\(\_\_\_\_\)

\subsection{题目5}
6．当 \(t\) 满足条件 \(\_\_\_\_\)时，二次型 \(f\left(x_{1}, x_{2}, x_{3}\right)=x_{1}^{2}+2 x_{2}^{2}+3 x_{3}^{2}+2 x_{1} x_{2}-2 x_{1} x_{3}+2 t x_{2} x_{3}\) 是正定的．

\subsection{题目6}
7．设 \(A=\left(\begin{array}{lll}0 & 1 & 0 \\ 0 & 0 & 1 \\ 0 & 0 & 0\end{array}\right)\) ，则与 \(A\) 乘积可交换的三阶方阵组成的线性空间的维数是 \(\_\_\_\_\)．\\
\(\_\_\_\_\)

\subsection{题目7}
8．已知 \(\sigma\) 为 \(n\) 维欧氏空间 \(V\) 上的线性变换，且 \(\sigma^{2}=0\) ，则 \(\sigma\) 的秩至多是 \(\_\_\_\_\)．

\subsection{题目8}
9．设 \(A=\left(\begin{array}{cc}2 & 1 \\ -1 & 0\end{array}\right)\) ，则 \(A\) 的2024次方幂是 \(\_\_\_\_\)．\\
\(\_\_\_\_\)

\subsection{题目9}
10．已知 \(V\) 是实数域上所有 \(n\) 元有序数组构成的欧氏空间，两个向量的内积定义为对应分量的乘积的和，则 \(V\) 上的柯西－布尼亚科夫斯基不等式的一般形式是 \(\_\_\_\_\)。

二．（10 分）已知 \(a_{1}, a_{2}, \cdots, a_{n-1}\) 是 \(n-1\) 个两两不同的数，请将 \(f(x)\) 分解成不可约一次多项式的乘积，并求出 \(f(x)\) 的根，其中

\[
f(x)=\left|\begin{array}{ccccc}
1 & x & x^{2} & \cdots & x^{n-1} \\
1 & a_{1} & a_{1}^{2} & \cdots & a_{1}^{n-1} \\
1 & a_{2} & a_{2}^{2} & \cdots & a_{2}^{n-1} \\
\vdots & \vdots & \vdots & & \vdots \\
1 & a_{n-1} & a_{n-1}^{2} & \cdots & a_{n-1}^{n-1}
\end{array}\right|
\]

三．（10 分）设关于字母 \(x\) 的 \(n\) 次首 1 多项式 \(f(x)\) 的一阶导数 \(\frac{\mathrm{d} f}{\mathrm{~d} x}\) 只有非零根，且 \(\frac{\mathrm{d} f}{\mathrm{~d} x}\) 整除 \(n+1\) 次多项式 \(x f(x)\) ，证明：\(f(x)\) 是某个一次多项式的 \(n\) 次方幂．

四．（10分）设 \(A=\left(a_{i j}\right), B=\left(b_{i j}\right)\) 是两个 \(n\) 阶方阵，证明：矩阵方程 \(A X=B\) 有解的充要条件是\(n+1\) 个矩阵 \(A, A_{1}, A_{2}, \cdots, A_{n}\) 有相同的秩，其中

\[
A_{k}=\left(\begin{array}{ccccc}
a_{11} & a_{12} & \cdots & a_{1 n} & b_{1 k} \\
a_{21} & a_{22} & \cdots & a_{2 n} & b_{2 k} \\
\vdots & \vdots & & \vdots & \vdots \\
a_{n 1} & a_{n 2} & \cdots & a_{n n} & b_{n k}
\end{array}\right), k=1,2, \cdots, n .
\]

五．（15 分）设实二次型 \(f(X)=X^{\mathrm{T}} A X\) 满足：\(X^{\mathrm{T}} A X=0\) 当且仅当 \(X=0\) ，证明：该二次型或者为正定二次型，或者为负定二次型。

六。（10 分）设 \(G\) 是数域 \(F\) 上一些 \(n\) 阶可逆方阵的集合，且满足：对 \(G\) 中任意不同的矩阵 \(A, B\) ，都有 \(|A+B|=0\) 。设 \(X=\left(x_{i j}\right)\) 是由 \(n^{2}\) 个字母 \(x_{i j}\) 组成的 \(n\) 阶方阵，对每一个 \(A \in G\) ，定义多项式

\[
f_{A}(X)=|X+A|
\]

（1）（5 分）证明：集合 \(\left\{f_{A}(X) \mid A \in G\right\}\) 的任意有限个子集都是线性无关的．\\
（2）（5 分）证明：集合 \(G\) 是有限的．

七．（15 分）设矩阵 \(A=\alpha \beta^{\mathrm{T}}\) ，其中 \(\alpha, \beta \in \mathbb{R}^{n}(n \geq 2), \alpha \neq 0, \beta \neq 0\) ．证明：\\
原题第（1）问：\(\alpha\) 以及与 \(\beta\) 正交的非零向量 \(\xi\) 都是矩阵 \(A\) 的特征向量．\\

\subsection{题目15}
（15 分）设矩阵 \(A=\alpha \beta^{\mathrm{T}}\) ，其中 \(\alpha, \beta \in \mathbb{R}^{n}(n \geq 2), \alpha \neq 0, \beta \neq 0\) ．证明：\\
原题第（2）问：若 \(\alpha\) 与 \(\beta\) 正交，则矩阵 \(A\) 可对角化．\\

\subsection{题目15}
（15 分）设矩阵 \(A=\alpha \beta^{\mathrm{T}}\) ，其中 \(\alpha, \beta \in \mathbb{R}^{n}(n \geq 2), \alpha \neq 0, \beta \neq 0\) ．证明：\\
原题第（3）问：若 \(\alpha\) 与 \(\beta\) 正交，则矩阵 \(A\) 不可对角化．

八．（15 分）设 \(\alpha\) 是 \(n(n \geq 1)\) 维欧氏空间 \(V\) 中的单位向量，定义 \(V\) 上的线性变换

\[
\sigma(\beta)=\beta-2(\alpha, \beta) \alpha
\]

（1）（5 分）证明：线性变换 \(\sigma\) 是正交变换（称这种变换为 \(V\) 上的镜面反射）．\\
（2）（5 分）证明：线性变换 \(\sigma\) 是第二类的（即行列式为－1）。\\
（3）（5 分）若上述 \(V\) 上的线性变换 \(\tau\) 以1为特征值，且属于特征值1的特征值子空间的维数是 \(n-1\)，证明：\(\tau\) 是 \(V\) 上的镜面反射．

九．（15 分）幂等方阵指的是平方不变的方阵，幂零方阵指的是某次方幂等于零的方阵。设 \(A\) 是复数域上的 \(n\) 阶方阵，证明：\\
（1）（7 分）\(A\) 可以表示为一个可逆方阵与一个幂等方阵之积。\\
（2）（8 分）\(A\) 可以表示为一个可对角化方阵和一个幂零方阵之和。

十．（10 分）设 \(\sigma\) 是 \(n\) 维线性空间上的可逆线性变换．\\
（1）（5 分）证明：\(\sigma\) 的逆变换可以表示成 \(\sigma\) 的多项式的形式。\\
（2）（5 分）设 \(f(x)\) 是 \(\sigma\) 的特征多项式，且多项式 \(f(x)\) 与 \(g(x)\) 是互素的，证明：\(g(\sigma)\) 是可逆线性变换。
\end{document}
